% 102-09(102쪽) — 삼차함수 y=f(x) 의 그래프(x 축과 -3, 0, 2 에서 만난다). 스캔 화질이 낮아 TikZ 로 다시 그림.
% 원문에 있는 것만 옮긴다: 곡선 · 라벨 y=f(x), -3, O, 2, x, y. y 축 눈금은 원문에 없으므로 넣지 않는다.
% 곡선은 발문의 조건 f(-3)=f(0)=f(2)=0, f(1)=-4 를 그대로 만족하는 모양으로 그렸다(원문 그림과 같은 개형).
\begin{tikzpicture}[x=0.30cm, y=0.25cm, line width=0.6pt]
  \draw[->, >=stealth, line width=0.7pt] (-4.60,0) -- (3.55,0) node[below=1pt, font=\small] {$x$};
  \draw[->, >=stealth, line width=0.7pt] (0,-4.60) -- (0,10.80) node[left=1pt, font=\small] {$y$};
  \draw[line width=0.6pt] plot[smooth, tension=0.6] coordinates {
    (-3.26,-4.46) (-3.15,-2.43) (-3.00,0) (-2.85,2.07) (-2.70,3.81) (-2.50,5.63)
    (-2.30,6.92) (-2.10,7.75) (-1.90,8.15) (-1.75,8.20) (-1.60,8.06) (-1.40,7.62)
    (-1.20,6.91) (-1.00,6.00) (-0.80,4.93) (-0.60,3.74) (-0.40,2.50) (-0.20,1.23)
    (0,0) (0.20,-1.15) (0.40,-2.18) (0.60,-3.02) (0.80,-3.65) (1.00,-4.00)
    (1.12,-4.06) (1.30,-3.91) (1.50,-3.38) (1.70,-2.40) (1.85,-1.35) (2.00,0)
    (2.15,1.66) (2.30,3.66) (2.45,6.01) (2.55,7.78) (2.64,9.53) };
  \node[below=1pt, font=\small] at (-3,0) {$-3$};
  \node[below left=-1pt and -1pt, font=\small] at (0,0) {$\pt{O}$};
  \node[below=1pt, font=\small] at (2,0) {$2$};
  \node[anchor=west, font=\small] at (0.45,9.95) {$y=f(x)$};
\end{tikzpicture}
